% Sección de John Silva: AMQP/RabbitMQ y análisis comparativo, aplicada al SSDC.
% La propuesta es una evolución futura: el SSDC actual es un monolito y NO usa un broker.
\begingroup
\section{AMQP y análisis comparativo}

\begin{frame}{AMQP y RabbitMQ: exchange, binding y cola}
{\small El productor publica en un \textbf{exchange}; los \textbf{bindings} lo conectan con \textbf{colas} y los consumidores confirman (\textit{ack}).}
\begin{center}
\resizebox{0.9\textwidth}{!}{
\begin{tikzpicture}[x=1cm,y=1cm]
  \node[caja,text width=2.5cm,minimum height=1.3cm] (pub) at (0,0)
    {\textbf{Servicio de prepostulaciones}\\\scriptsize productor};
  \node[evento,text width=2.5cm,minimum height=1.3cm] (ex) at (4.3,0)
    {\textbf{Exchange topic}\\\scriptsize\texttt{ssdc.eventos}};
  \node[caja,text width=2.5cm] (q1) at (8.8,1.6)  {\scriptsize\texttt{q.correo}};
  \node[caja,text width=2.5cm] (q2) at (8.8,0)    {\scriptsize\texttt{q.notificaciones}};
  \node[caja,text width=2.5cm] (q3) at (8.8,-1.6) {\scriptsize\texttt{q.auditoria}};
  \node[caja,text width=2.3cm] (c1) at (12.6,1.6)  {\scriptsize Servicio de correo};
  \node[caja,text width=2.3cm] (c2) at (12.6,0)    {\scriptsize Servicio de notificaciones};
  \node[caja,text width=2.3cm] (c3) at (12.6,-1.6) {\scriptsize Servicio de auditoría};
  \draw[flecha] (pub) -- node[above,font=\scriptsize]{publica} (ex);
  \foreach \q in {q1,q2,q3}{\draw[flecha] (ex.east) -- ++(0.6,0) |- (\q.west);}
  \foreach \q/\c in {q1/c1,q2/c2,q3/c3}{\draw[flecha] (\q) -- (\c);}
  \node[font=\tiny\ttfamily,text=gris] at (6.65,2.35) {binding: prepostulacion.aprobada};
  \node[font=\scriptsize,text=gris] at (10.7,-2.45) {consumo con \textit{ack}};
\end{tikzpicture}}
\end{center}
{\small\textbf{Evento de ejemplo:} \texttt{PrepostulacionAprobada} (hoy, dentro del backend y sin broker).}\par\vspace{1mm}
{\scriptsize\color{gris}Matiz: exchange y binding son del modelo 0-9-1; AMQP 1.0 (OASIS) no los define.}
\fuentes{\cite{oasis2012,rabbitmq_concepts}}
\end{frame}

\begin{frame}{MQTT vs JMS vs AMQP: ¿cuál encaja en el SSDC?}
\small
\begin{center}
\renewcommand{\arraystretch}{1.25}
\resizebox{\textwidth}{!}{%
\begin{tabular}{@{}p{2.6cm}p{4.2cm}p{4.2cm}p{4.6cm}@{}}
\toprule
 & \textbf{MQTT} & \textbf{JMS} & \textbf{AMQP} \\
\midrule
Naturaleza & Protocolo ligero pub/sub & API Java; el proveedor define el protocolo & Protocolo de red estándar \\
Modelo & Tópicos & Colas y tópicos & Exchanges, bindings y colas \\
Entrega & QoS 0, 1 y 2 & Ack y mensajes persistentes & Ack, colas durables, confirmaciones \\
Encaje en el SSDC & Pensado para dispositivos con poco ancho de banda; el SSDC no los tiene & Integra bien con Spring, pero ata el diseño a Java & Independiente del lenguaje; permite enrutar a varios servicios \\
\bottomrule
\end{tabular}}
\end{center}
\begin{block}{Tolerancia a fallos del broker}
\small Si el broker cae, los servicios no se comunican: hace falta colas durables, \textit{ack} manual, cola de mensajes fallidos y un broker replicado.
\end{block}
\fuentes{\cite{derlatka2023,maharjan2023,bezerra2024,gultunca2018}}
\end{frame}

\begin{frame}{Recomendación para el SSDC y compromisos}
\footnotesize
\begin{columns}[T,onlytextwidth]
\column{0.52\textwidth}
\textbf{AMQP con RabbitMQ para eventos asíncronos}
\begin{itemize}
  \item Un exchange \texttt{topic} para todos los eventos.
  \item Colas durables y \textit{ack} manual.
  \item Cola de fallidos y reintentos acotados.
  \item Consumidores idempotentes (\texttt{eventoId}).
\end{itemize}
\column{0.44\textwidth}
\textbf{Compromisos}
\begin{itemize}
  \item Un componente más que operar y monitorear.
  \item Consistencia eventual: el correo llega después de aprobar.
  \item Entrega al menos una vez: pueden llegar duplicados.
\end{itemize}
\end{columns}
\begin{block}{Antes de adoptarlo}
\footnotesize Medir con la carga real del SSDC: los estudios comparativos no son mediciones de nuestro sistema \cite{maharjan2023}.
\end{block}
{\scriptsize\color{gris}Propuesta de evolución: el SSDC actual no usa broker.}
\fuentes{\cite{rabbitmq_concepts,maharjan2023}}
\end{frame}

\begin{frame}{Incidente real: Google Cloud, 2 de junio de 2019}
\small
\begin{center}
\resizebox{0.97\textwidth}{!}{
\begin{tikzpicture}[x=1cm,y=1cm]
  \node[caja,text width=3.2cm] (a) at (0,0)    {\textbf{Detonante}\\\scriptsize mantenimiento en una sola ubicación};
  \node[evento,text width=3.6cm] (b) at (4.6,0) {\textbf{Causa raíz}\\\scriptsize error de software en la automatización: detuvo clústeres de lugares distintos a la vez};
  \node[caja,text width=3.4cm] (c) at (9.4,0)  {\textbf{Defensa que falló}\\\scriptsize modo \textit{fail static} duró solo minutos; se retiraron rutas BGP};
  \node[caja,text width=3.0cm] (d) at (13.8,0) {\textbf{Efecto}\\\scriptsize pérdida de capacidad de red en seis regiones};
  \draw[flecha] (a) -- (b); \draw[flecha] (b) -- (c); \draw[flecha] (c) -- (d);
\end{tikzpicture}}
\end{center}
\begin{itemize}
  \item Duración: entre 3 h 19 min y 4 h 25 min según la región.
  \item Herramientas de depuración compitieron con la red congestionada.
\end{itemize}
\begin{block}{Qué aprendemos para el SSDC con broker}
\footnotesize Un cambio de configuración o de mantenimiento no debe aplicarse a todos los nodos del broker a la vez: se prueba primero en uno y se puede volver atrás.
\end{block}
\fuentes{\cite{googlecloud19009}}
\end{frame}
\endgroup
